% Book p. 50 (PDF p. 51)
\section*{Modelling with Algebra}
\subsection*{Application of Algebra}

\subsubsection*{Introduction}

Studying sequences and inequalities is important because arithmetic and
geometric sequences provide statistical understanding into central tendencies,
growth rates, etc. Also, quadratic inequalities help in decision-making processes
where certain conditions or constraints must be met. For example, in business
and economics, they are used to model profit maximisation or loss minimisation
under specific constraints. In physics and engineering, quadratic inequalities arise
when determining safe operating conditions, or analysing forces. In architecture
and design, they help define areas, volumes or structural integrity within specific
boundaries. This section introduces you to the concept of the sum of sequences,
convergent and divergent series, recursive sequences, arithmetic and geometric
sequences. You will also learn how to solve quadratic inequalities, systems of
quadratic inequalities and real-life problems involving quadratic inequalities.

\begin{keyideas}
\begin{itemize}
  \item A \textbf{sequence} is a set of numbers in a given order.
  \item Each of the numbers in a sequence is called a term.
  \item Maximum values are the largest possible values of variables within a
  system that satisfy a set of constraints, each represented by an inequality.
  \item Minimum values are the least possible values of variables within a system
  that satisfy a set of constraints, each represented by an inequality.
  \item When finding the solution to a quadratic inequality, it is advisable to start
  by sketching the associated quadratic graph.
\end{itemize}
\end{keyideas}
